\documentclass{article}

\usepackage[utf8]{inputenc}
\usepackage[danish]{babel}
\usepackage{lipsum}
\usepackage{enumitem}
\setlist{label=(\alph*), topsep=0cm}
\usepackage{amsmath, amssymb, amsthm}
\makeatletter
\renewcommand*\env@matrix[1][*\c@MaxMatrixCols c]{%
  \hskip -\arraycolsep
  \let\@ifnextchar\new@ifnextchar
  \array{#1}}
\makeatother
\usepackage{graphicx}
\usepackage{hyperref}
\usepackage[bb=boondox]{mathalfa}
\usepackage{xcolor}
\usepackage{chngcntr}
\usepackage{mathtools}
\counterwithin{equation}{section}
\renewcommand{\qedsymbol}{$\blacksquare$}

\newcommand{\upint}[2]{
  \overline{\int_{#1}^{#2}}
}
\newcommand{\loint}[2]{
  \underline{\int_{#1}^{#2}}
}

% Margen
\usepackage[left=2.5cm, right=2.5cm, top=3cm, bottom=3cm]{geometry}

% Eksempel Enviroment
\newcounter{eksempel}[section]
\newenvironment{eksempel}[1][]{\refstepcounter{eksempel}\par\medskip
   \noindent \textbf{Eksempel~\theeksempel. #1} \rmfamily}{\medskip}

% Spacing omkring align enviroment
% \usepackage{etoolbox}
% \newcommand{\zerodisplayskips}{%
%   \setlength{\abovedisplayskip}{-8pt}%       % Spacing over Align
%   \setlength{\belowdisplayskip}{2pt}%        % Spacing under Align
%   \setlength{\abovedisplayshortskip}{0pt}%   % Spacing over Equation
%   \setlength{\belowdisplayshortskip}{0pt}}%  % Spacing under Equation
% \appto{\normalsize}{\zerodisplayskips}%
% \appto{\small}{\zerodisplayskips}%
% \appto{\footnotesize}{\zerodisplayskips}%

%	¤¤ Afsnitsformatering ¤¤ %
\setlength{\parindent}{0mm}           		% Stoerrelse af indryk
\setlength{\parskip}{3mm}          			% Afstand mellem afsnit ved brug af double Enter
\linespread{1,1}							% Linie afstand

\usepackage{mdframed}
\mdtheorem[%
linecolor=black,
  linewidth=.5pt,
  backgroundcolor=gray!10,
  frametitlerule=true,
  frametitlefont=\bfseries,
  frametitlebackgroundcolor=gray!25,
  innertopmargin=\topskip]{theorem}{Theorem}[section]

\mdtheorem[%
linecolor=black,
  linewidth=.5pt,
  backgroundcolor=gray!10,
  frametitlerule=true,
  frametitlefont=\bfseries,
  frametitlebackgroundcolor=gray!25,
  innertopmargin=\topskip]{corollary}[theorem]{Korollar}

\mdtheorem[%
linecolor=black,
  linewidth=.5pt,
  backgroundcolor=gray!10,
  frametitlerule=true,
  frametitlefont=\bfseries,
  frametitlebackgroundcolor=gray!25,
  innertopmargin=\topskip]{lemma}[theorem]{Lemma}

\mdtheorem[%
linecolor=black,
  linewidth=.5pt,
  backgroundcolor=gray!10,
  frametitlerule=true,
  frametitlefont=\bfseries,
  frametitlebackgroundcolor=gray!25,
  innertopmargin=\topskip]{proposition}[theorem]{Proposition}

\mdtheorem[%
linecolor=black,
  linewidth=.5pt,
  backgroundcolor=gray!10,
  frametitlerule=true,
  frametitlefont=\bfseries,
  frametitlebackgroundcolor=gray!25,
  innertopmargin=\topskip]{definition}[theorem]{Definition}

%   Stort spørgsmålstegn kommando
\newcommand{\?}{
\begin{center}
    \huge{\textbf{?}}
\end{center}
}

% Math in section
\newcommand{\mbf}[1]{$\mathbf{#1}$}

%kasse til sætninger mv.
\newenvironment{boks}[1]{\begin{mdframed}\textbf{#1.}}{\end{mdframed}}

\DeclareMathOperator{\poi}{Poi}
\DeclareMathOperator{\unif}{unif}
\DeclareMathOperator{\var}{Var}

\begin{document}

\begin{center}
\thispagestyle{empty}
\parskip=14pt%
\vspace*{3\parskip}%

\begin{figure}[h]
\centering
    \includegraphics[scale=2]{Figures/AAU_logo/AAU_LOGO_RGB.png}
\end{figure}

\vspace{-8pt}
\rule{16cm}{1pt}
\vspace{5pt}

\Huge{\textbf{Eksamensnoter}}\\
\vspace{8pt}
\Large{\textbf{4. semester:} \textit{Analyse 2}}\\
\vspace{50pt}

\Large{\textbf{Kasper Rosenkrands}}
\end{center}

\newpage


\pagenumbering{roman}
\renewcommand{\baselinestretch}{0.95}\normalsize
\tableofcontents
\renewcommand{\baselinestretch}{1.0}\normalsize
\newpage

\pagenumbering{arabic}

\section{Differentiable funktioner af flere variable}
% !TEX root = ../main.tex
\subsection{ABC-kriteriet}
\begin{boks}{Sætning 9.26 (ABC-kriteriet)}
  Lad $D \subseteq \mathbb{R}^2$ og lad $f \ : \ D \rightarrow \mathbb{R}$ være en funktion. Antag, at $z_0 = (x_0, y_0) \in D$ er et kritisk punkt for $f$, dvs. $z_0$ ligger i det indre af $D$ og $\nabla f(z_0)$. Antag yderligere, at de partielt afledte af $f$ af anden orden eksisterer og er kontinuerte i en åben kugle $B_r (z_0)$. Sæt
  \begin{align*}
    A = \frac{\partial^2 f}{\partial x^2}(z_0), \quad B = \frac{\partial^2 f}{\partial x \partial y}(z_0), \quad C = \frac{\partial^2 f}{\partial y^2}(z_0).
  \end{align*}
  \begin{enumerate}
    \item Hvis $B^2 > AC$, har $f$ et saddelpunkt i $z_0$, og
    \item hvis $B^2 < AC$, har $f$ et strengt lokalt ekstremum i $z_0$.

    Det lokale ekstremum i (b) er et strengt lokalt minimum, hvis $A > 0$ og $C > 0$, og det er et strengt lokalt maksimum, hvis $A < 0$ og $C < 0$ ($A$ og $C$ har samme fortegn, fordi $AC > 0$).
    \item Hvis $B^2 = AC$, giver $ABC$-kriteriet ingen oplysning om $f$'s opførsel i nærheden af $z_0$.
  \end{enumerate}
\end{boks}
\begin{proof}
  (i) Vi skal bevise resultaterne ved at se på $f$'s opførsel på linjer gennem $z_0$, så lad $v = (h,k)$ være en enhedsvektor i $\mathbb{R^2}$ og lad $g$ være funktionen defineret ved
  \begin{align*}
    g(t) = f(z_0 + tv) = f(x_0 + th, y_0 + tk)
  \end{align*}
  og med delmængden
  \begin{align*}
    D_g = \{t \in \mathbb{R} \ | \ z_0 + tv \in D\}
  \end{align*}
  af $\mathbb{R}$ som sin definitionsmængde.

  Da $B_r(z_0) \subseteq D$ og $\| v \| = 1$, er $]-r,r[ \subseteq D_g$.

  Da $f$'s partielt afledte af anden orden eksisterer og er kontinuerte i $B_r(z_0)$, følger det af sætning 9.12, at de partielt afledte af første orden er differentiable i hele denne kugle.

  Så er $f$'s partielt afledte af første orden også kontinuerte i $B_r(z_0)$, og igen får vi af Sætning 9.12, at funktionen $f$ selv er differentiabel i $B_r(z_0)$.

  Kædereglen giver nu, at $g$ er differentiabel i intervallet $]-r,r[$ med diferentialkvotienten
  \begin{align}\label{9.35}
    g'(t) = \frac{\partial f}{\partial x} (x_0 + th, y_0 + tk)h + \frac{\partial f}{\partial y}(x_0 + th, y_0 + tk)k.
  \end{align}

  Da $\frac{\partial f}{\partial x}$ og $\frac{\partial f}{\partial y}$ også er differentiable i $B_r(z_0)$, kan Kædereglen bruges på udtrykket for $g'(t)$ på højre side af denne ligning, og vi får, at $g'$ er differentiabel i intervallet $]-r,r[$ med differentialkvotienten
  \begin{align*}
    g''(t) = &\left[ \frac{\partial^2 f}{\partial x^2}(x_0 + th, y_0 + tk)h + \frac{\partial^2 f}{\partial y \partial x}(x_0 + th, y_0 + tk)k \right] h\\
    &+ \left[ \frac{\partial^2 f}{\partial x \partial y}(x_0 + th, y_0 + tk)h + \frac{\partial^2 f}{\partial y^2}(x_0 + th, y_0 + tk) k \right] k.
  \end{align*}
  Med henblik på at skrive dette lidt mere overskueligt indfører vi betegnelserne
  \begin{align}
    A(x,y) &= \frac{\partial^2 f}{\partial x^2}(x,y)\nonumber\\
    B(x,y) &= \frac{\partial^2 f}{\partial x \partial y}(x,y), \quad \text{og} \label{9.36}\\
    C(x,y) &= \frac{\partial^2 f}{\partial y^2}(x,y).\nonumber
  \end{align}
  Idet vi husker, fra sætningen om blandede partielt afledete og fordi de afledte er kontiuerte
  \begin{align*}
    \frac{\partial^2 f}{\partial x \partial y}(x,y) = \frac{\partial^2 f}{\partial y \partial x}(x,y) = B(x,y)
  \end{align*}
  ifølge sætning 9.18, har vi så
  \begin{align}\label{9.37}
    g''(t) = A(z_0 + tv)h^2 + 2B(z_0 + tv)hk + C(z_0 + tv)k^2.
  \end{align}
  Resten af beviset for sætningen deler vi op i de to tilfælde, som der er skelnet imellem i sætningen.

  (ii) Lad os først betragte tilfældet
  \textit{(a)} $B^2 > AC$.
  Ifølge Lemma 9.25 eksisterer der vektorer $u_1 = (s_1, t_1)$ og $u_2 = (s_2, t_2)$ således, at
  \begin{align*}
    As_1^2 + 2Bs_1t_1 + Ct_1^2 > 0 \quad \text{og} \quad As_2^2 + 2Bs_2t_2 + Ct_2^2 < 0.
  \end{align*}

  Lad $v_1 = (h_1, k_1)$ og $v_2 = (h_2, k_2)$ være de enhedsvektorer, vi får ved at normere $u_1$ og $u_2$, altså
  \begin{align*}
    v_1 = \frac{u_1}{\| u_1 \|} \quad \text{og} \quad v_2 = \frac{u_2}{\| u_2 \|},
  \end{align*}
  og lad $g_1$ og $g_2$ være de to funktioner
  \begin{align*}
    g_1(t) &= f(z_0 + tv_1) = f(x_0 + th_1, y_0 + tk_1), \quad \text{0}\\
    g_2(t) &= f(z_0 + tv_2) = f(x_0 + th_2, y_0 + tk_2).
  \end{align*}
  Ifølge (i) er de begge definerede og to gange differentiable i intervallet $]-r,r[$. Da $\nabla f(x_0, y_0) = 0$, får vi af \eqref{9.35}, at
  \begin{align*}
    g_1'(0) = \frac{\partial f}{\partial x}(x_0, y_0)h_1 + \frac{\partial f}{\partial y}(x_0, y_0)k_1 = 0 \quad \text{og} \quad g_2'(0) = 0.
  \end{align*}
  I udtrykket
  \begin{align*}
    g_1''(t) = A(z_0 + tv_1)h_1^2 + 2B(z_0 + tv_1)h_1k_1 + C(z_0 + tv_1)k_1^2
  \end{align*}
  for differentialkvotienten $g_1''$ af anden orden indgår funktionen
  \begin{align*}
    (x,y) \mapsto A(x,y)h_1^2 + 2B(z_0 + tv_1)h_1k_1 + C(x,y)k_1^2.
  \end{align*}

  Den er kontinuert på $B_r(z_0)$, og da
  \begin{align*}
    A(x_0, y_0)&h_1^2 + 2B(x_0, y_0)h_1k_1 + C(x_0, y_0)k_1^2\\
    &= Ah_1^2 + 2Bh_1k_1 + C k_1^2 = \frac{As_1^2 + 2Bs_1t_1 + Ct_1^2}{\| u_1 \|^2}>0,
  \end{align*}
  findes der et $\rho > 0$ således, at $\rho < r$ og
  \begin{align*}
    A(x,y)h_1^2 + 2B(x,y)h_1k_1 + C(x,y)k_1^2 > 0
  \end{align*}
  for alle $(x,y) \in B_\rho (z_0)$.

  Altså er $g''(t) > 0$ for alle $|t| < \rho$, og af Korrolar 7.18 følger så, at funktionen $g_1$ på intervallet $]-\rho, \rho[$ har \textit{strengt} globalt minimum i 0, dvs.
  \begin{align*}
    f(z_0 + tv_1) = g_1(t) > g_1(0) = f(z_0)
  \end{align*}
  for $0 < |t| < \rho$.

  Heraf ses, at $f$ \textit{ikke} har lokalt maksimum i $z_0$.

  Betragtning afa funktionen $g_2$ viser, at $f$ heller ikke har lokalt minimum i $z_0$, men det vil jo sige at $z_0$ er et saddelpunkt for $f$.

  (iii) Lad os dernæst betragte tilfældet
  \textit{(b)} $B^2 < AC$
  Lad os tillige antage, at $A > 0$ og $C > 0$ (tilfældet $A < 0$ og $C < 0$ følger heraf ved at skifte fortegn på $f$).

  Ifølge antagelserne er de tre funktioner
  \begin{align*}
    A(x,y), \quad B(x,y) \quad \text{og} \quad C(x,y)
  \end{align*}
  i \eqref{9.36} kontinuerte på $B_r(z_0)$, og så er funktionen
  \begin{align*}
    (x, y) \mapsto (B(x,y))^2 - A(x,y)C(x,y)
  \end{align*}
  det også.

  Da
  \begin{align*}
    (B(x_0, y_0))^2 - A(x_0, y_0)C(x_0, y_0) = B^2 - AC < 0
  \end{align*}
  og
  \begin{align*}
    A(x_0, y_0) = A > 0,
  \end{align*}
  findes der et $\rho > 0$ således, at $\rho < r$ og
  \begin{align*}
    (B(x,y))^2 - A(x_0, y_0)C(x_0, y_0) = B^2 - AC < 0
  \end{align*}
  og
  \begin{align*}
    (B(x,y))^2 - A(x,y)C(x,y) < 0 \quad \text{og} \quad A(x,y) > 0
  \end{align*}
  for alle $(x,y) \in B_\rho(z_0)$.

  Af Lemma 9.25 får vi, at
  \begin{align*}
    A(x,y)h^2 + 2B(x,y)hk + C(x,y)k^2 > 0
  \end{align*}
  for alle $(x,y) \in B_\rho(z_0)$ og for enhver enhedsvektor $v = (h,k) \in \mathbb(R)^2$.

  Lad da $v = (h,k)$ være en vilkårlig enhedsvektor i $\mathbb{R}^2$ og betragt som i (i) funktionen
  \begin{align*}
    g(t) = f(z_0 + tv) = f(x_0 + th, y_0 + tk).
  \end{align*}

  Som i (ii) ovenfor ser vi, at $g'(0) = 0$ og at $g''(t) > 0$ for alle $|t| < \rho$. Af Korrolar 7.18 følger så, at betragtet på intervallet $]-\rho, \rho[$ har $g$ strengt globalt minimum i 0, dvs.
  \begin{align*}
    f(z_0 + tv) = g(t) > g(0) = f(z_0)
  \end{align*}
  for $0 < |t| < \rho$.

  Lad nu $z = (x,y)$ være et vilkårligt fra $z_0 = (x_0, y_0)$ forskelligt punkt i $B_\rho(z_0)$. Et sådant punkt kan skrives på formen
  \begin{align*}
    z = z_0 + tv, \quad \text{hvor} \quad \| v \| = 1 \quad \text{og} \quad 0 < t < \rho,
  \end{align*}
  og da det foregående gælderr for alle $v$, ser vi, at
  \begin{align*}
    f(x,y) > f(z_0) \quad \text{for alle} \quad (x,y) \in B_\rho(z_0),
  \end{align*}
  og altså, at $f$ har strengt lokalt minimum i $z_0$.
\end{proof}

\subsection{Differentiabilitetssætningen i det genrelle tilfælde}
\begin{boks}{Sætning 9.12}
  Lad $f \ : \ U \rightarrow \mathbb{R}$ være en funktion og lad $x_0 \in U$. Antag, at de partielt afledte af $f$ eksisterer i alle punkter i en åben kugle $B_r(x_0) \subseteq U$, og at de derved fremkomne funktioner
  \begin{align*}
    \frac{\partial f}{\partial x_j} \ : \ B_r(x_0) \rightarrow \mathbb{R}, \
    j = 1,2,\ldots,n,
  \end{align*}
  alle er kontinuerte i $x_0$. Så er $f$ differentiabel i $x_0$.
\end{boks}
\begin{proof}
  Ideerne i beviset er de samme som i beviset for tilfældet $n = 2$, så vi gør ikke så meget ud af detaljerne. Lad $x_0$ have koordinaterne $x_0 = (a_1, a_2, \ldots, a_n)$ og lad $x = (x_1, x_2, \ldots, x_n)$ være et vilkårligt punkt i $b_r(x_0)$. Sæt
  \begin{align*}
    y^k = (x_1, x_2, \ldots, x_k, a_{k + 1}, a_{k + 2}, \ldots, a_n)
    \quad \text{for} \quad k = 0, 1, 2, \ldots, n.
  \end{align*}
  Bemærk, at $y^0 = x_0$, $y^n = x$ og $y^k \in B_r(x_0)$ for alle $k$, samt at
  \begin{align*}
    f(x) - f(x_0) = \sum_{k = 1}^n f(y^k) - f(y^{k - 1}).
  \end{align*}
  Idet alle koordinater for $y^k$ og $y^{k - 1}$ er de samme med undtagelse af den $k$'te, fås af Middelværdisætningen, at
  \begin{align*}
    f(y^k) - f(y^{k - 1}) = \frac{\partial f}{\partial x_k}(z_k)(x_k - a_k),
  \end{align*}
  hvor $z_k$ er et punkt på linjestykket, der forbinder $y^{k - 1}$ med $y^k$. Vi har altså
  \begin{align*}
    f(x) - f(x_0) &= \sum_{k = 1}^n \frac{\partial f}{\partial x_k}(z_k)(x_k - a_k) \\
    &= \langle \nabla f(x_0), x - x_0 \rangle\\
    &\quad + \sum_{k = 1}^n \left( \frac{\partial f}{\partial x_k}(z_k) - \frac{\partial f}{\partial x_k}(x_0) \right) (x_k - a_k).
   \end{align*}
   Heraf følger resultatet, hvis vi siger, at der eksisterer en o-funktion $\varphi$ således, at
   \begin{align*}
     \sum_{k = 1}^n \left( \frac{\partial f}{\partial x_k}(z_k) - \frac{\partial f}{\partial x_k}(x_0) \right) (x_k - a_k) =
     \varphi(x - x_0) \| x -x_0 \|.
   \end{align*}
   Da denne ligning er opfyldt for $x = x_0$, hvis $\varphi(0) = 0$, skal vi blot bevise, at
   \begin{align*}
     \sum_{k = 1}^n \left( \frac{\partial f}{\partial x_k}(z_k) - \frac{\partial f}{\partial x_k}(x_0) \right) \frac{x_k - a_k}{\| x - x_0 \|} \rightarrow 0 \quad \text{for} \quad x \rightarrow x_0.
   \end{align*}
   Lad $\varepsilon > 0$ være givet. For ethvert $k$ findes et $\delta_k > 0$ således, at der for alle $x \in B_r(x_0)$ gælder
   \begin{align*}
     \| x - x_0 \| \delta_k \quad \Rightarrow \quad \left| \frac{\partial f}{\partial x_k}(x) - \frac{\partial f}{\partial x_k}(x_0) \right| < \frac{\varepsilon}{n}.
   \end{align*}
   Vælg $\delta$ som det mindste af tallene $r$ og $\delta_1, \delta_2, \ldots, \delta_n.$

   Da $|x_k - a_k| \leq \| x - x_0 \|$ for alle $k = 1,2, \ldots, n,$ får vi for alle $x$ i den udprikkede kugle $\dot{B}_\delta (x_0)$
   \begin{align*}
     \bigg| \sum_{k = 1}^n &\bigg( \frac{\partial f}{\partial x_k}(z_k) - \frac{\partial f}{\partial x_k}(x_0) \bigg) \frac{x_k - a_k}{\| x - x_0 \|} \bigg| \\
     &\leq \sum_{k = 1}^n \left|\frac{\partial f}{\partial x_k}(z_k) - \frac{\partial f}{\partial x_k}(x_0)\right| < n\frac{\varepsilon}{n} = \varepsilon,
   \end{align*}
   hvormed sætningen er bevist.
\end{proof}

\newpage

\section{Taylors formel}
% !TEX root = ../main.tex
Måske vil jeg have Taylors Sætning for funktioner af én variabel med i stedet for Sætning 2.2.

\subsection{Taylors Sætning for funktioner af flere variable}
Før vi formulerer og beviser Taylors Sætning for funktioner af flere variable, vil vi først introducere en meget nyttig og - anvendt notation.
\begin{boks}{Definition 2.1 (Multi-indeks-notation)}
  Et \textit{multi-indeks} er en $n$-tupel af ikke-negative heltal $\alpha \in \mathbb{N}_0^n$. For to multi-indices $\alpha = (\alpha_1, \alpha_2, \ldots, \alpha_n)$, $\beta = (\beta_1, \beta_2, \ldots, \beta_n) \in \mathbb{N}_0^n$ og et $x = (x_1, x_2, \ldots, x_n) \in \mathbb{R}^n$ defineres:
  \begin{enumerate}[label = \arabic*.]
    \item Sum/differens: $\alpha \pm \beta = (\alpha_1 \pm \beta_1, \alpha_2 \pm \beta_2, \ldots, \alpha_n \pm \beta_n).$
    \item  Absolutværdi: $|\alpha| = \sum_{i = 1}^n \alpha_i.$
    \item Fakultet: $\alpha! = \prod_{i = 1}^n \alpha_i!.$
    \item Potens: $x^\alpha = \prod_{i = 1}^n x_i^{\alpha_i}.$
    \item Partielt afledet af højere orden: $\partial^\alpha = \prod_{i = 1}^n \partial_i^{\alpha_i}$ hvor $\partial_i^{\alpha_i} = \frac{\partial^{\alpha_i}}{\partial x_i^{\alpha_i}}.$
  \end{enumerate}
\end{boks}

\begin{boks}{Sætning 2.2 (Taylors Sætning med restled)}
  Lad $k \in \mathbb{N}$, $x,x_0 \in A \subset \mathbb{R}^n$, $x \neq x_0$ og $f \ : \ A \rightarrow \mathbb{R}$ opfylde, at $\partial^\alpha f$ eksisterer og er differentiabel på linjestykket $L = \{(1 - t)x_0 + tx \in \mathbb{R}^n \ | \ t \in [0,1]\}$ mellem $x_0$ og $x$ for $|\alpha| \leq k$.
  Så eksisterer et $y \in \{(1 - t)x_0 + tx \in \mathbb{R}^n \ | \ t \in (0,1)\}$ så
  \begin{align*}
    f(x) = \sum_{|\alpha| \leq k} \frac{\partial^\alpha f(x_0)}{\alpha!}(x - x_0)^\alpha + \sum_{|\alpha| = k + 1}\frac{\partial^\alpha f(y)}{\alpha!}(x - x_0)^\alpha.
  \end{align*}
\end{boks}
\begin{proof}
  Da $\partial^\alpha f$ er antaget differentiabel på $L$ for alle $|\alpha| \leq k$, så må $L \subset A$ og $x_0, x \in A$ være indre punkter, og vi kan finde et $r > 0$ så $F \ : \ (-r, 1 + r) \rightarrow \mathbb{R}$ givet ved $F(t) = f((1-t)x_0 + tx)$ er veldefineret og $k + 1$ gange differentiabel på $[0,1]$. Dermed kan vi anvende Sætning 1.1 på $F$ med $x_0 = 0$ og $x = 1$:
  \begin{align*}
    F(1) = \sum_{j = 0}^k \frac{F^{(j)}(0)}{j!} + \frac{F^{(k + 1)}(c)}{(k + 1)!},
  \end{align*}
  hvor $c \in (0,1)$. Kædereglen giver nu:
  \begin{align*}
    F^{(j)}(t) = \sum_{|\alpha| = j} \frac{j!}{\alpha!}(\partial^\alpha f)((1-t)x_0 + tx)(x - x_0)^\alpha,
  \end{align*}
  se Opgave 1, hvoraf resultatet følger med $y = (1 - c)x_0 + cx$.
\end{proof}

\subsection{Lokale ekstrema for funktioner af flere variable}
\begin{boks}{Definition 4.4 (Hesse-matrix)}
  Lad $x \in A \subset \mathbb{R}^n$, $A$ åben og $f \in C^2(A)$. Så kaldes matricen $H(x)$ givet ved
  \begin{align*}
    h_{ij}(x) = \frac{\partial^2 f(x)}{\partial x_i \partial x_j}, \quad
    H(x) = (h_{ij}(x))_{i,j = 1}^n
  \end{align*}
  for \textit{Hesse-matricen} for $f$ i $x$.
\end{boks}

\begin{boks}{Lemma 4.5}
  Lad $x,x_0 \in A \subset \mathbb{R}^n$ opfylde at $A$ er åben og at $L = \{(1-t)x_0 + tx \in \mathbb{R}^n \ | \ t \in (0,1)\}$ så
  \begin{align*}
    f(x) = f(x_0) + \langle \nabla f(x_0), x - x_0 \rangle + \frac{1}{2} \langle x - x_0, H(y)(x - x_0) \rangle.
  \end{align*}
\end{boks}
\begin{proof}
  Beviset er blot Taylors sætning med restled for $k = 1$:
  \begin{align}
    f(x) &= \sum_{| \alpha | \leq 1} \frac{\partial^\alpha f(x_0)}{\alpha!} (x - x_0)^\alpha +
    \sum_{| \alpha | = 3} \frac{\partial^\alpha f(y)}{\alpha!} (x - x_0)^\alpha \nonumber\\
    &= \frac{\partial^0 f(x_0)}{0!}(x - x_0)^0 +
    \frac{\partial^1 f(x_0)}{1!}(x - x_0)^1 +
    \frac{1}{2}\frac{\partial^2 f(y)}{\partial x_1^2}(x - x_0)^2 +
    \frac{1}{2}\frac{\partial^2 f(y)}{\partial x_2^2}(x - x_0)^2\nonumber\\
    &\qquad+
    \frac{1}{2}\frac{\partial^2f(y)}{\partial x_1 \partial x_2}(x - x_0) +
    \frac{1}{2}\frac{\partial^2 f(y)}{\partial x_1 \partial x_2}(x - x_0) +
    \frac{1}{2}\frac{\partial^2 f(y)}{\partial x_2 \partial x_1}(x - x_0)\nonumber\\
    &= f(x_0) + f'(x_0)(x - x_0) +
    1/2 \left( \frac{\partial^2f(y)}{\partial x_1 \partial x_2}(x - x_0)^2 +
    \frac{\partial^2 f(y)}{\partial x_2^2}(x - x_0)^2 + 2\frac{\partial^2 f(y)}{\partial x_1 \partial x_2} \right)\label{taylor_omskrivning}
  \end{align}
  Her kommer nogle mellem trin, de følgende matricer bruges til forståelse
  \begin{align*}
    H(y) =
    \begin{bmatrix}
      \frac{\partial^2 f(y)}{\partial x_1^2} &
      \frac{\partial^2 f(y)}{\partial x_2 \partial x_1} \\
      \frac{\partial^2 f(y)}{\partial x_1 \partial x_2} &
      \frac{\partial^2 f(y)}{\partial x_2^2}
    \end{bmatrix}, \qquad
    H(y)(x - x_0) =
    \begin{bmatrix}
      \frac{\partial^2 f(y)}{\partial x_1^2}(x - x_0) &
      \frac{\partial^2 f(y)}{\partial x_2 \partial x_1}(x - x_0) \\
      \frac{\partial^2 f(y)}{\partial x_1 \partial x_2}(x - x_0) &
      \frac{\partial^2 f(y)}{\partial x_2^2}(x - x_0)
    \end{bmatrix}.
  \end{align*}
  Ved brug af det ovenstående ser vi at
  \begin{align*}
    &\frac{1}{2} \langle x - x_0, H(y)(x - x_0) \rangle =\\
    &\qquad \frac{1}{2}
    \bigg( \frac{\partial^2 f(y)}{\partial x_1^2}(x- x_0)^2 +
    \frac{\partial^2 f(y)}{\partial x_2 \partial x_1}(x- x_0)^2 +
    \frac{\partial^2 f(y)}{\partial x_1 \partial x_2}(x- x_0)^2 +
    \frac{\partial^2 f(y)}{\partial x_2^2}(x- x_0)^2 \bigg)
  \end{align*}
  Nu kan vi omskrive \eqref{taylor_omskrivning} til det følgende, ved brug af det ovenstående
  \begin{align*}
    f(x_0) + \langle \nabla f(x_0), x - x_0 \rangle + \frac{1}{2}\langle x - x_0, H(y)(x - x_0)\rangle
  \end{align*}
  hvilket fuldender beviset.
\end{proof}

\begin{boks}{Definition 4.6}
  Antag at $A \subset \mathbb{R}^n$ er åben, at $f \in C^2(A)$, at $x_0 \in A$ er et kritisk punkt for $f$ og at alle Hesse-matricen $H(x_0)$'s egenværdier er positive.
  Så er $x_0$ et lokalt minimumspunkt for $f$. Hvis alle $H(x_0)$'s egenværdier omvendt er negative, er $x_0$ et lokalt maksimumspunkt for $f$.
\end{boks}
\begin{proof}
  Antag uden tab af generalitet at alle egeværdier for $H(x_0)$ er positive.
  Da $H(x_0)$ er reel og symmestrisk er den selvadjungeret og spektralsætningen giver, at $H(x_0)$ har $n$ reelle, ortonormale egenvektorer $e_i$, $i = 1, \ldots, n$ med tilhørende egenværdier $\lambda_i$. Vi kan derfor for $z \in \mathbb{R}^n$ skrive
  \begin{align*}
    z = \sum_{i = 1}^n \langle z, e_i \rangle e_i
  \end{align*}
  og dermed
  \begin{align*}
    H(x_0)z = \sum_{i = 1}^n \langle z, e_i \rangle \lambda_i e_i
  \end{align*}
  eller
  \begin{align*}
    \langle z, H(x_0)z \rangle = \sum_{i = 1}^n \lambda_i \langle z, e_i \rangle^2.
  \end{align*}
  Lad $m = \min \{ \lambda_i \ | \ i \in \{ 1, \ldots, n \}\} > 0$. Så er
  \begin{align*}
    \langle z, H(x_0)z \rangle = \sum_{i = 1}^n \lambda_i \langle z, e_i \rangle^2 \geq \sum_{i = 1}^n m \langle z, e_i \rangle^2 = m \| z \|^2.
  \end{align*}
  Sæt $A(x) = H(y) - H(x_0)$, som afhænger af $x$ gennem $y \in \{ (1 - t)x_0 + tx \in \mathbb{R}^n \ | \ t \in (0,1) \}$.
  Da $f \in C^2(A)$ er $\| A(x) \|_{HS}$ kontinuert og vi kan finde $r > 0$ så $x \in B_r(x_0) \backslash \{ x_0 \}$ medfører, at $\| A(x) \|_{HS} \leq \frac{1}{2}m$. Anvender vi nu Lemma 4.5 på $x \in B_r(x_0)$ fås
  \begin{align*}
    f(x) &= f(x_0) + \langle \nabla f(x_0), x - x_0 \rangle + \frac{1}{2}\langle x - x_0, H(y)(x - x_0)\rangle \\
    &= f(x_0) + 0 + \frac{1}{2}\Big( \langle x - x_0, H(x_0)(x - x_0) \rangle + \langle x - x_0, A(x)(x - x_0) \rangle \Big)\\
    &\geq f(x_0) + \frac{1}{2}\| x - x_0 \|^2 (m - \frac{1}{2}m) =
    f(x_0) + \frac{m\| x - x_0 \|^2}{4},
  \end{align*}
  hvor vi igen brugte Cauchy-Schwartz.
\end{proof}

\newpage

\section{Sætningen om inverse funktioner}
% !TEX root = ../main.tex

\newpage

\section{Sætningen om implicit givne funktioner}
% !TEX root = ../main.tex
% \subsection{Lemma 7.2}
% \begin{boks}{Lemma 7.2}
%   \begin{itemize}[label = $\ast$]
%     \item Lad $\mathcal{O} \ \dot{\subset} \ \mathbb{R}^d$, $K \subset \mathcal{O}$ hvor $K$ er konveks og kompakt
%     \item $\varphi \in C^1(\mathcal{O})$, dvs. at de partielt afledte eksisterer og er kontinuerte.
%     \item Definér
%     $$ \| \partial_j \varphi \|_\infty \ := \ \sup_{x \in \mathcal{O}} | \partial_j \varphi(x) | < \infty$$
%   \end{itemize}
%   Så er $\varphi$ lipschitz-kontinuert på $K$, dvs.
%   $$ | \varphi(x) - \varphi(x') | \leq L \| x - x' \|, \quad \forall x,x' \in K$$
%   hvor
%   $$ L \ := \ \sqrt{\sum_{j = 1}^d \| \partial_j \varphi} \|^2_\infty $$
% \end{boks}
% At $K$ er en konveks mængde betyder at
% $$ (1 - t)x' + tx \in K, \quad \forall x,x' \in K, \ \forall t \in [0,1] $$
% \begin{proof}
%   Lad $x,x' \in K$.
%   Definér $f \ : \ [0,1] \rightarrow \mathbb{R}$ som
%   $$ f(t) = \varphi\Big( (1 - t)x' + tx \Big), \quad 0 \leq t \leq 1 $$
%   Vi ved $\varphi \in C^1(\mathcal{O})$, så da $K \subset \mathcal{O}$ og konveks er $f$ differentiabel og kædereglen fortæller os at
%   $$ f'(t) = \langle \nabla \varphi\big( (1 - t)x' + tx \big), x - x' \rangle $$
%   Nu kigger på vi venstre siden af den ulighed vi gerne vil vise
%   \begin{align*}
%     | \varphi(x) - \varphi(x') | =
%     | f(1) - f(0) | =
%     \left| \int_0^1 f'(t) \right| \leq
%     \int_0^1 |f'(t)| dt
%   \end{align*}
%   vi kan nu indsætte det kendte udtryk for $f'(t)$
%   \begin{align*}
%     \int_0^1 | \langle \nabla \varphi\big( (1 - t)x' + tx \big), x - x' \rangle| dt &\leq \int_0^1 \| \nabla \varphi \big( (1 - t)x' + tx \big) \| \cdot \| x - x' \| dt \\
%     &= \int_0^1 \sqrt{\sum_{j = 1}^d | \partial_j \varphi\big( (1 - t)x' + tx \big) |^2 } \cdot \| x - x' \| dt \\
%     &\leq \int_0^1 \sqrt{\sum_{j = 1}^d \| \partial_j \varphi \|^2 } \cdot \| x - x' \| dt \\
%     &= L \| x - x' \| \int_0^1 1 dt
%   \end{align*}
% \end{proof}

\begin{boks}{Sætning 7.4}
  \begin{itemize}[label = $\ast$]
    \item $U \ \dot{\subseteq} \ \mathbb{R}^{d = m + n}$
    \item $h \in C^1(U, \mathbb{R}^m)$
    \item Antag at der eksisterer et punkt $a = [u_a, w_a]$ så $h(a) = 0$
    \item og at $[D_u h(a)]$ er invertibel
  \end{itemize}

  Så findes $f \ : \ E \ \dot{\subseteq} \in \mathbb{R}^n \rightarrow \mathbb{R}^m$ så $f(w_a) = u_a$, $h([f(w), w]) = 0 \ \forall w \in E$.
  % 
  % Og $f \in C^1(E, \mathbb{R}^m)$ med $[Df(w)] = -[D_u h([f(w), w])]^{-1}[D_w h([f(w), w])], \forall w \in E$ hvor $[D_u h([f(w), w])]$ er invertibel.
\end{boks}

\begin{proof}
  For ethvert $w \in P_n(\varepsilon)$, definer da funktionen $F_w \ : \ \overline{Q_m(\varepsilon)} \rightarrow \mathbb{R}^m$, givet ved
  \begin{align*}
    F_w(u) \ := \ u - [D_u h(a)]^{-1} h([u,w])
  \end{align*}
  Da skal der fines et $C > 0$ og et $\varepsilon < \frac{r}{\sqrt{2}}$, sådan at der for ethvert $\varepsilon \leq \varepsilon_1$ og for alle $\|w - w_a\| < \frac{\varepsilon}{2C}$ gælder følgende
  \begin{enumerate}[label = \arabic*.]
    \item $F_w\bigg( \overline{Q_m(\varepsilon)} \bigg) \subseteq \overline{Q_m(\varepsilon)}$
    \item $F_w : \overline{Q_m(\varepsilon)} \rightarrow \overline{Q_m(\varepsilon)}$ er en kontraktion.
  \end{enumerate}

  Jeg starter med at vise at funktionen $F_w$ er en kontraktion.

  Det vil med andre ord sige at der eksisterer et $0 < \varepsilon_1 < \frac{r}{\sqrt{2}}$ småt nok til, at der for ethvert $\varepsilon \leq \varepsilon_1$ og $w \in P_n(\varepsilon)$ gælder
  \begin{align}\label{enkelt}
    \| F_w(u) - F_w(u') \| \leq \frac{1}{2} \| u - u' \|, \quad \forall u,u' \in \overline{Q_m(\varepsilon)}
  \end{align}
  Vi kan omskrive $F_w$ til
  \begin{align*}
    F_w(u) = -[D_uh(a)]^{-1}(h([u,w]) - [D_uh(a)]u)
  \end{align*}
  Definer nu
  \begin{align*}
    g_w(u) = h([u,w]) - [D_uh(a)]u
  \end{align*}
  dvs.
  \begin{align*}
    F_w(u) = -[D_uh(a)]^{-1}g_w(u)
  \end{align*}
  iflg. Lemma 7.1
  \begin{align*}
    \| F_w(u) - F_w(u') \| \leq \|[D_uh(a)]^{-1} \|_{HS} \| g_w(u) - g_w(u') \|
  \end{align*}

  Fordi $\overline{Q_w(\varepsilon)}$ er kompakt og konveks kan vi bruge Lemma 7.3 til at vise det følgende
  \begin{align*}
    \| g_w(u) - g_w(u') \| \leq \| \Delta g_w \|_{\infty, \overline{Q_w(\varepsilon)}} \| u - u' \|, \quad \forall u, u' \in \overline{Q_w(\varepsilon)}, \ w \in P_n(\varepsilon)
  \end{align*}
  Da $g_w$'s partielt afledte er kontinuerte, eksisterer der et $\varepsilon_1 > 0$ så
  \begin{align*}
    \| \Delta g_w \|_{\infty, \overline{Q_w(\varepsilon_1)}} \leq \frac{1}{2\| [D_uh(a)]^{-1} \|_{HS}}, \quad \forall w \in P_n(\varepsilon_1)
  \end{align*}
  indsæt i \eqref{enkelt} og indse
  \begin{align}\label{tredobbelt}
    \| F_w(u) - F_w(u') \| \leq \frac{1}{2} \| u - u' \|
  \end{align}
  Vi kan altså konkludere at $F_w$ opfylder kontraktionsuligheden.

  Nu viser vi at $\overline{Q_w(\varepsilon)}$ er invariant under $F_w$.
  Med andre ord $F_w : \overline{Q_w(\varepsilon_1)} \rightarrow \overline{Q_w(\varepsilon_1)}$ for $w \in P_n(\frac{\varepsilon}{2C})$, C konstant og $\varepsilon$ tilpas lille. Ifølge Lemma 7.1 er
  \begin{align} \label{dobbelt}
    \| F_w(u) - F_{w'}(u) \| \leq \| [D_uh(a)]^{-1} \|_{HS}
    \| h([u,w]) - h([u,w']) \|, \quad \forall u \in \overline{Q_w(\varepsilon)} \ \text{og} \ w,w' \in \mathbb{R}^m
  \end{align}
  Vi vil nu vurdere det sidste udtryk ovenfor, til dette benyttes Lemma 7.3, benytter mængden $\Omega = \overline{Q_w(\varepsilon)} \times \overline{P_n(\varepsilon_1)}$ som er kompakt og konveks.
  Vi får uligheden
  \begin{align*}
      \| h([u,w]) - h([u,w']) \| \leq L \| [u,w] - [u,w'] \|
  \end{align*}
  Vi indsætter i \eqref{dobbelt}
  \begin{align*}
    \| F_w(u) - F_{w'}(u) \| &\leq \| [D_uh(a)]^{-1} \|_{HS} L\| [u,w] - [u,w'] \| \\
    &\leq \| [D_uh(a)]^{-1} \|_{HS} L\| w - w' \|, \quad \forall u \in \overline{Q_w(\varepsilon)} \ \text{og} \ w,w' \in \overline{P_n(\varepsilon_1)}
  \end{align*}
  Definer nu $C := \| [D_uh(a)]^{-1} \|_{HS}L$
  \begin{align*}
    \| F_w(u) - F_{w'}(u) \| \leq C\| w - w' \|, \quad \forall w,w' \in P_n(\varepsilon_1), \ \forall u \in \overline{Q_w(\varepsilon_1)}
  \end{align*}
  Vi viser nu $F_w(u) \in \overline{Q_w(\varepsilon)}$
  \begin{align*}
    \| F_w(u) - u_a \| = \| F_w(u) - F_{w_a}(u_a) \|
    &\leq \| F_w(u) - F_w(u_a) \| + \| F_w(u_a) - F_{w_a}(u_a) \| \\
    &\leq \frac{1}{2} \| u - u_a \| + C \| w - w_a \| \\
    &< \frac{\varepsilon}{2} + C\frac{\varepsilon}{2C} = \varepsilon
   \end{align*}
   Dermed er $F_w \in \overline{Q_m(\varepsilon)}$, og dermed er $F_w$ en kontraktion.

   $(\overline{Q_m(\varepsilon)}, d_{\mathbb{R}^m})$ opfylder kriterierne for et fuldstændigt metrisk rum. Dermed vil alle Cauchy følger i $\overline{Q_m(\varepsilon)}$ konvergerer med grænse i $\overline{Q_m(\varepsilon)}$, som følge af følge karakterisation af lukkede mængder.

   Derfor kan vi nu tillade os at kalde $\overline{Q_m(\varepsilon)}$ for et Banach rum og derfor eksisterer et entydigt fikspunkt $u_w$ for $F_w$
   $$F(u_w) = u_w$$
   Vi ser at
   $$F(u_w) = u_w - [D_uh(a)]^{-1}h(u_w, w) = u_w$$
   dermed må
   $$[D_uh(a)]^{-1}h(u_w, w) = 0 $$
   $$h(u_w, w) = 0$$
   Vi definerer $f:P_n(\frac{\varepsilon_1}{2C}) \rightarrow \mathbb{R}^m$
   $$f(w) = u_w$$
   Vi har altså at
   \begin{align*}
     f(w_a) = F_{w_a}(u_a) &= u_a - [D_uh(a)]^{-1}h(w_a, u_a)\\
     &= u_a - 0 \\
     &= u_a
   \end{align*}
   hvor $E := P_n(\frac{\varepsilon_1}{2C})$.
\end{proof}

% \textbf{Nu vil den tredje del af beviset for Sætningen om implicit givne funktioner blive gennemgået}
%
% \begin{proof}
%   Lad $w, w' \in P_n\left(\frac{\varepsilon}{2C}\right)$. Vi ved at $f(w), f(w') \in \overline{Q_m(\varepsilon)}$.
%
%   Sæt
%   \begin{align*}
%   x = [f(w), w], \
%   x' = [f(w'), w']
%   \in \overline{Q_m(\varepsilon)} \times  P_n\left(\frac{\varepsilon}{2C}\right)
%   \end{align*}
%   Vi ved at $h$ er differentiabel i $x'$ dvs.
%   \begin{align}\label{imp_kors}
%     h(x) = h(x') + [Dh(x')](x - x') + \varphi_{x'}(x - x')\|x - x'\|
%   \end{align}
%   Betragt funktionen $$ x \mapsto \det[D_uh(x)] $$
%   $[D_uh(a)]$ er antaget invertibel derfor må $\det[D_uh(a)] \neq 0$.
%   $h$ er desuden kontinuert differentiabel hvilket mefører at
%   $$ \exists \ r > 0 \ \forall x \in B_r(a) \ : \ \det[D_uh(x)] \neq 0. $$
%   Der kan derfor vælges $\varepsilon \leq r$ så $[D_uh(x)]$ er defineret for $x \in E = B_\varepsilon (a)$
%
%   Vi kan nu omskrive \eqref{imp_kors} med vores $x = [f(w), w]$ og $x' = [f(w'), w']$, samt opdelingen af Jacobi matricen
%   $$[Dh(x)] = \Big[ [D_uh(x)], [D_wh(x)] \Big] $$
%
%   \begin{align*}
%     h([f(w), w]) &= h([f(w'), w']) +
%      \Big[ [D_uh([f(w'), w'])], [D_wh([f(w'), w'])] \Big]([f(w), w] - [f(w'), w']) + \\
%      &\qquad\varphi_{x'}([f(w), w] - [f(w'), w']) \| [f(w), w] - [f(w'), w'] \| \\
%     &= h([f(w'), w']) +
%     [D_uh([f(w'), w'])](f(w) - f(w')) +
%     [D_wh(f(w'), w'])](w - w') +  \\
%     &\qquad\varphi_{x'}\Big( [f(w) - f(w')], [w - w'] \Big) \| [f(w) - f(w')], [w - w'] \|
%   \end{align*}
%   Husk at $h(f(w), w) = h(f(w'), w') = 0$, og at $\det[D_uh(f(w'), w')] \neq 0 \ \forall w \in P_n\left( \frac{\varepsilon}{2C} \right)$
%   Dermed fås det følgende ved at gange med den inverse og indse at nogle af ledene er nul
%   \begin{align*}
%     f(w) - f(w') &= -[D_uh(f(w'), w')]^{-1}[D_wh(f(w'), w')](w - w') -[D_uh(f(w'), w')]^{-1} \varphi_{x'}(\ldots)\| \ldots \|
%   \end{align*}
%   Det er nok at vise at $$ f(w) - f(w') = L(w - w') + \varphi(w - w') \| w - w'\| $$ for $L$ lineær og $\varphi$ en lille-o funktion.
%   Vi ser straks at matrix-vektor produktet er en lineær afbildning.
%   Nu mangler vi blot at vise o-funktionen, bemærk at
%   \begin{align*}
%     \|[f(w) - f(w'), w - w'] \|
%   \end{align*}
%   er begrænset.
%
%   Vi har altså nogle konstanter og en lille-o funktion.
%   $f$ er derfor differentiabel i $w'$ og dermed på $P_n\left( \frac{\varepsilon}{2C} \right)$.
% \end{proof}

\newpage

\section{Banachs fikspunktsætning}
% !TEX root = ../main.tex

\newpage

\section{Eksistens og entydighed af løsninger til ODE'er}
% !TEX root = ../main.tex
\begin{boks}{Sætning 6.3}
  Antag $U \ \dot{\subseteq} \ \mathbb{R}^d$, $I \ \dot{\subseteq} \ \mathbb{R}$. Antag $y_0 \in U$ og $r_0, \delta_0 > 0$ så $\overline{B_{r_0}(y_0)} \subset U$, $[t_0 - \delta_0, t_0 + \delta_0] \subset I$.
  Betragt $f \ : \ I \times U \rightarrow \mathbb{R}^d$ kontinuert hvor om der gælder for $L > 0$.
  \begin{align}\label{enkelt_kors}
    \| f(t,x) - f(t,y)\| \leq L \| x - y \|, \quad \forall t \in [t_0 - \delta_0, t_0 + \delta_0], \quad \forall x,y \in B_{r_0}(y_0)
  \end{align}
  så for IVP:
  \begin{align}\label{trippel_kors}
    y'(t) = f(t, y(t)), \quad y(t_0) = y_0.
  \end{align}
  Definer $\delta_1 := \min \{\delta_0, r_0/M, 1/L \}$. Da eksisterer en entydig løsning
  \begin{align*}
    y \ : \ (t - \delta_1, t_0 + \delta_1) \rightarrow \overline{B_{r_0}(y_0)}
  \end{align*}
\end{boks}
\begin{proof}
  Tag $ 0 < \delta < \delta_1$, og definer $K \ := \ [t_0 - \delta, t_0 + \delta] \subset \mathbb{R}$ enhver kontinuert funktion $\phi \ : \ K \rightarrow \mathbb{R}^d$ er født begrænset og fordi $\mathbb{R}^d$ er et Banach rum siger proposition 6.2 at rummet $(C(K;\mathbb{R}^d), d_\infty)$ er et fuldstændigt metrisk rum.

  Definer
  \begin{align}\label{dobbelt_kors}
    X \ := \ \{ g \in C(K;\mathbb{R}^d) \ : \ g(t) \in \overline{B_{r_0}(y_0)}, \forall t \in K \}
  \end{align}
  \begin{boks}{Lemma 6.4}
    Det metriske rum $(X, d_\infty)$ er fuldstændigt.
  \end{boks}
  \begin{proof}
    Betragt en Cauchy følge $\{ g_n \}_{n\geq 1} \subset X$. Fordi $(C(K;\mathbb{R}^d), d_\infty)$ er fuldstændigt har vi en $g \in (C(K;\mathbb{R}^d)$ hvor $\lim_{n \rightarrow \infty} d_\infty(g_n, g) = 0$. Vi har med andre ord for ethvert $t\in K$
    \begin{align*}
      g(t) = \lim_{n \rightarrow \infty} g_n(t), \quad \lim_{n \rightarrow \infty} \| g_n(t) - g(t) \| = 0.
    \end{align*}
    Fordi følgen er i $X$ er det følgende opfyldt $\| g_n(t) - y_0 \| \leq r_0 \quad \forall t,n$ hvilket medfører at
    \begin{align*}
      \| g(t) - y_0 \| = \lim_{n \rightarrow \infty} \| g_n(t) - y_0 \| \leq r_0
    \end{align*}
    og $g$ tilhører dermed $X$ hvilket konkluderer beviset.
  \end{proof}
  \begin{boks}{Lemma 6.5}
    Definer $F \ : \ X \rightarrow C(K; \mathbb{R}^d)$
    \begin{align*}
      [F(g)](t) \ := \ y_0 + \int_{t_0}^t f(s,g(s))ds, \quad \forall t \in K
    \end{align*}
    hvor $f$ opfylder \eqref{enkelt_kors}. Hvorfor (i) $vm_f \subseteq X$ og (ii) $F \ : \ X \rightarrow X$ er en kontraktion.
  \end{boks}
  \begin{proof}
    \textbf{udeladt}
  \end{proof}
  \textbf{Forsættelse af beviset for Sætning 6.3}

  $F$ kontraktion altså siger Banach at der eksisterer et entydig kontinuert funktion $ y \ : \ K \rightarrow \overline{B_{r_0}(y_0)}$ så
  \begin{align*}
    y(t) = [F(y)](t) = y_0 + \int_{t_0}^tf(s,y(s))ds, \quad t\in [t_0 - \delta, t_0 + \delta]
  \end{align*}
  dette kan differentieres som følge af Analysens Fundamental Sætning
  \begin{align*}
    \frac{\partial}{\partial t}\left( \int_{t_0}^t f(s,y(s))ds \right) =
    f(t, y(t)) = y'(t)
  \end{align*}
  hvilket viser \eqref{trippel_kors}.
\end{proof}

\newpage

\end{document}
